\documentclass[a4paper,14pt]{extarticle}

\usepackage{geometry}
\geometry{a4paper,tmargin=2cm,bmargin=2cm,lmargin=3cm,rmargin=1cm}
\usepackage[T2A]{fontenc}
\usepackage[utf8]{inputenc}
\usepackage[english,main=russian]{babel}
\usepackage{indentfirst}
\usepackage{graphicx}
\graphicspath{{./images/}}
\usepackage{setspace}
\linespread{1.3}
\usepackage{tempora}
\usepackage{amsmath}
\usepackage{booktabs}
\usepackage{array}
\usepackage{caption}
\usepackage{hyperref}
\hypersetup{unicode=true,hidelinks}
\usepackage{pgfplots}
\pgfplotsset{compat=1.18}
\usepackage{listings}

\sloppy
\clubpenalty=10000
\widowpenalty=10000

\lstdefinelanguage{Go}{
  morekeywords={break,default,func,interface,select,case,defer,go,map,struct,
    chan,else,goto,package,switch,const,fallthrough,if,range,type,continue,
    for,import,return,var},
  morecomment=[l]{//},
  morecomment=[s]{/*}{*/},
  morestring=[b]",
}
\lstset{
  language=Go,
  basicstyle=\ttfamily\footnotesize,
  breaklines=true,
  breakatwhitespace=true,
  columns=fullflexible,
  frame=single,
  numbers=left,
  numbersep=5pt,
  xleftmargin=18pt,
  framexleftmargin=18pt,
  showstringspaces=false,
  tabsize=2,
  captionpos=t,
  inputencoding=utf8,
  keepspaces=true
}
\renewcommand{\lstlistingname}{Листинг}

\begin{document}

\begin{titlepage}
\thispagestyle{empty}

\vspace*{-60pt}

\hspace*{-82pt}
\begin{minipage}{0.3\textwidth}
\centering
\IfFileExists{images/emblem.png}{%
  \includegraphics[width=0.9\textwidth]{emblem.png}%
}{%
  \IfFileExists{images/emblem.jpg}{%
    \includegraphics[width=0.9\textwidth]{emblem.jpg}%
  }{}%
}
\end{minipage}
\begin{minipage}{0.67\textwidth}
\small
\textbf{
\vspace*{-0.7ex}
\centerline{Министерство науки и высшего образования Российской Федерации}
\vspace*{-0.7ex}
\centerline{Федеральное государственное бюджетное образовательное учреждение}
\vspace*{-0.7ex}
\centerline{высшего образования}
\vspace*{-0.7ex}
\centerline{<<Московский государственный технический университет}
\vspace*{-0.7ex}
\centerline{имени Н. Э. Баумана}
\vspace*{-0.7ex}
\centerline{(национальный исследовательский университет)>>}
\vspace*{-0.7ex}
\centerline{(МГТУ им. Н. Э. Баумана)}
}
\end{minipage}

\vspace{-25pt}
\hspace{-35pt}\rule{\textwidth}{2.3pt}

\vspace*{-20.3pt}
\hspace{-35pt}\rule{\textwidth}{0.4pt}

\vspace{1.5ex}
\hspace{-35pt}\noindent\small
ФАКУЛЬТЕТ\hspace{80pt} <<Информатика и системы управления>>

\vspace*{-16pt}
\hspace{47pt}\rule{0.83\textwidth}{0.4pt}

\vspace{0.5ex}
\hspace{-35pt}\noindent\small
КАФЕДРА\hspace{50pt} <<Теоретическая информатика и компьютерные технологии>>

\vspace*{-16pt}
\hspace{30pt}\rule{0.866\textwidth}{0.4pt}

\vspace{11em}

\begin{center}
\Large \textbf{Лабораторная работа №1} \\
\large \textbf{по курсу <<Параллельные и распределённые системы>>} \\
\large <<Распараллеливание алгоритма вычисления произведения двух матриц>>
\end{center}

\normalsize

\vspace{8em}

\begin{flushright}
Студент группы ИУ9-51Б \\
Шалимов Д. И. \\
\vspace{2ex}
Преподаватель Вишняков И. Э.
\end{flushright}

\bigskip
\vfill

\begin{center}
\textsl{Москва 2026}
\end{center}
\end{titlepage}

\setcounter{page}{2}
\tableofcontents
\newpage

\section{Цель работы}

Цель работы --- реализовать стандартный алгоритм умножения двух квадратных
матриц без использования библиотечных функций матричного умножения, сравнить
обход результата по строкам и по столбцам, а также получить многопоточную
реализацию и оценить её ускорение.

\section{Постановка задачи}

Пусть даны две квадратные матрицы $A$ и $B$ порядка $n$. Необходимо вычислить
матрицу $C = AB$ по формуле
\[
    c_{ij} = \sum_{k=0}^{n-1} a_{ik}b_{kj}.
\]

Требовалось:
\begin{enumerate}
    \item реализовать последовательное умножение с обходом элементов $C$ по
    строкам и по столбцам;
    \item разбить матрицу $C$ на приблизительно равные прямоугольные блоки;
    \item поручить вычисление каждого блока отдельному worker;
    \item дождаться завершения всех worker и проверить полученную матрицу
    сравнением с последовательным результатом;
    \item измерить время при различном числе worker.
\end{enumerate}

\section{Реализация}

Программа написана на языке Go. Матрица хранится в одном непрерывном массиве
значений типа \texttt{int64}; элемент с индексами $(i, j)$ расположен по
индексу $i \cdot n + j$. Такой способ хранения уменьшает число выделений памяти
по сравнению с представлением в виде массива строк.

Последовательные функции \texttt{multiplyByRows} и
\texttt{multiplyByColumns} используют одну и ту же формулу умножения, но
меняют порядок обхода элементов результирующей матрицы. Время измеряется
функциями пакета \texttt{time}.

В параллельном варианте функция \texttt{chooseGrid} выбирает близкое к квадрату
разбиение числа worker на число блоков по строкам и столбцам. Затем
\texttt{splitRange} определяет границы блоков. Для каждого прямоугольного блока
матрицы $C$ запускается goroutine, которая вызывает \texttt{multiplyBlock}.
Записываемые областями $C$ не пересекаются, поэтому взаимное исключение не
требуется. Главная goroutine ожидает всех работников с помощью
\texttt{sync.WaitGroup}.

После вычисления каждого параллельного результата он сравнивается с
последовательным эталоном. Сравнение выполняется после остановки таймера и не
входит в измеряемое время.

Основная команда для параллельного измерения имеет вид:
\begin{lstlisting}[language={},caption={Команда запуска измерения}]
go run . -mode parallel -size 2000 -workers 1,2,4,8,16,32,64 -runs 1
\end{lstlisting}

Результат выполнения команды:
\begin{lstlisting}[language={},caption={Вывод программы для матриц 2000 x 2000}]
Matrix size: 2000 x 2000
Runs: 1
Logical CPUs: 16
Average sequential: 69.692656 s

Workers:  1 | Average: 121.451177 s | Speedup: 0.57x | Correct: true
Workers:  2 | Average: 65.202028 s | Speedup: 1.07x | Correct: true
Workers:  4 | Average: 33.891092 s | Speedup: 2.06x | Correct: true
Workers:  8 | Average: 19.425680 s | Speedup: 3.59x | Correct: true
Workers: 16 | Average: 11.251349 s | Speedup: 6.19x | Correct: true
Workers: 32 | Average: 18.142209 s | Speedup: 3.84x | Correct: true
Workers: 64 | Average: 18.133077 s | Speedup: 3.84x | Correct: true
\end{lstlisting}

Параметр \texttt{-runs} задаёт количество повторов. При значении больше одного
программа выводит среднее арифметическое измеренных времён.

\section{Условия эксперимента}

Измерения выполнены на следующем устройстве:
\begin{itemize}
    \item операционная система --- Windows 11 Enterprise 24H2;
    \item процессор --- Intel Core i5-13500H, базовая частота 2.60 ГГц;
    \item физических ядер --- 12;
    \item логических ядер --- 16;
    \item оперативная память --- 16 ГБ, 5200 MT/s;
    \item версия Go --- 1.27.0.
\end{itemize}

Перед началом вычислений программа устанавливает
\texttt{GOMAXPROCS} равным количеству доступных логических процессоров, то есть
16. Матрицы заполняются псевдослучайными числами от $-10$ до $10$ с
фиксированными начальными значениями генератора, поэтому входные данные
воспроизводимы.

\section{Результаты измерений}

\subsection{Серия измерений для матриц 1000 \texorpdfstring{$\times$}{x} 1000}

В этой серии для каждого варианта выполнено по десять запусков. В таблице
\ref{tab:sequential1000} приведены результаты последовательного режима.

\begin{table}[ht]
\centering
\caption{Последовательное умножение матриц 1000 \texorpdfstring{$\times$}{x} 1000}
\label{tab:sequential1000}
\begin{tabular}{lc}
\toprule
Порядок обхода $C$ & Среднее время, с \\
\midrule
По строкам & 2.928694 \\
По столбцам & 3.035877 \\
\bottomrule
\end{tabular}
\end{table}

В отдельном параллельном измерении среднее время последовательного эталона,
вычисленного в той же серии, составило 2.365027 с. Отличие от таблицы
\ref{tab:sequential1000} связано с тем, что это разные серии запусков; на
время влияют текущая загрузка и изменение частоты процессора.

\begin{table}[ht]
\centering
\caption{Параллельное умножение матриц 1000 \texorpdfstring{$\times$}{x} 1000, 10 запусков}
\label{tab:parallel1000}
\begin{tabular}{rrrr}
\toprule
Worker & Среднее время, с & Ускорение & Проверка \\
\midrule
1  & 2.476245 & 0.96 & true \\
2  & 1.267968 & 1.87 & true \\
4  & 0.714003 & 3.31 & true \\
8  & 0.505790 & 4.68 & true \\
16 & 0.503849 & 4.69 & true \\
32 & 0.510546 & 4.63 & true \\
64 & 0.529164 & 4.47 & true \\
\bottomrule
\end{tabular}
\end{table}

\subsection{Измерение для матриц 2000 \texorpdfstring{$\times$}{x} 2000}

Для более крупной матрицы выполнен один запуск каждого варианта
(\texttt{-runs 1}). Последовательный эталон занял 69.692656 с. Все результаты
параллельных вычислений совпали с ним.

\begin{table}[ht]
\centering
\caption{Параллельное умножение матриц 2000 \texorpdfstring{$\times$}{x} 2000}
\label{tab:parallel2000}
\begin{tabular}{rrrr}
\toprule
Worker & Время, с & Ускорение & Проверка \\
\midrule
1  & 121.451177 & 0.57 & true \\
2  & 65.202028  & 1.07 & true \\
4  & 33.891092  & 2.06 & true \\
8  & 19.425680  & 3.59 & true \\
16 & 11.251349  & 6.19 & true \\
32 & 18.142209  & 3.84 & true \\
64 & 18.133077  & 3.84 & true \\
\bottomrule
\end{tabular}
\end{table}

На рисунке \ref{fig:time2000} приведён график времени вычисления в зависимости
от числа worker для матриц порядка 2000. По оси абсцисс указано число worker,
по оси ординат --- время вычисления в секундах.

\begin{figure}[ht]
\centering
\begin{tikzpicture}
\begin{axis}[
    ybar,
    width=0.92\textwidth,
    height=8cm,
    xlabel={Число worker},
    ylabel={Время, с},
    symbolic x coords={1,2,4,8,16,32,64},
    xtick=data,
    ymin=0,
    ymax=130,
    grid=major,
    bar width=13pt,
    nodes near coords,
    every node near coord/.append style={font=\scriptsize, rotate=90, anchor=west}
]
\addplot coordinates {
    (1,121.451177)
    (2,65.202028)
    (4,33.891092)
    (8,19.425680)
    (16,11.251349)
    (32,18.142209)
    (64,18.133077)
};
\end{axis}
\end{tikzpicture}
\caption{Зависимость времени умножения матриц 2000 \texorpdfstring{$\times$}{x} 2000 от числа worker}
\label{fig:time2000}
\end{figure}

\section{Анализ результатов}

Увеличение числа worker от 1 до 16 уменьшает время вычисления. Для матриц
2000 \texorpdfstring{$\times$}{x} 2000 наибольшее ускорение равно 6.19 и
достигается при 16 worker. Это число совпадает с количеством логических ядер
процессора.

Один worker работает медленнее последовательного варианта: добавляются расходы
на создание goroutine, разбиение на блоки и ожидание в \texttt{WaitGroup}.
После 16 worker производительность уменьшается. При 32 и 64 worker несколько
задач конкурируют за одни и те же вычислительные ресурсы, а расходы планировщика
и влияние кэш-памяти становятся заметнее.

Windows позволяет изменять приоритеты процессов и потоков, однако в проведённых
измерениях приоритет не менялся: все запуски выполнялись с нормальным
приоритетом. Поэтому сравнение с высоким и низким приоритетами не приводится;
для такого сравнения необходимо выполнить отдельные серии запусков в одинаковых
условиях.

\section{Вывод}

Реализовано последовательное и параллельное умножение квадратных матриц на Go.
Корректность каждого параллельного результата подтверждена сравнением с
последовательным эталоном. Наиболее эффективным для проведённого измерения
оказалось использование 16 worker: для матриц 2000 \texorpdfstring{$\times$}{x}
2000 время уменьшилось с 69.692656 с до 11.251349 с, что соответствует
ускорению 6.19 раза. Увеличение числа worker сверх числа логических ядер не
привело к дополнительному ускорению.

\appendix
\newpage
\section{Листинги программы}

\begin{lstlisting}[caption={Файл 	exttt{main.go}}]
package main

import (
    "flag"
    "fmt"
    "os"
    "runtime"
    "strconv"
    "strings"
    "time"
)

func main() {
    mode := flag.String("mode", "parallel", "benchmark mode: sequential or parallel")
    size := flag.Int("size", 1000, "square matrix size")
    workerList := flag.String("workers", "1,2,4,8,16,32,64", "comma-separated worker counts")
    runs := flag.Int("runs", 10, "number of measured runs")
    flag.Parse()

    if *size <= 0 {
        fail("-size must be positive")
    }
    if *runs <= 0 {
        fail("-runs must be positive")
    }
    runtime.GOMAXPROCS(runtime.NumCPU())
    a := generateMatrix(*size, -10, 10, 1)
    b := generateMatrix(*size, -10, 10, 2)

    switch *mode {
    case "sequential":
        runSequential(a, b, *runs)
    case "parallel":
        workers, err := parseWorkers(*workerList)
        if err != nil {
            fail(err.Error())
        }
        runParallel(a, b, workers, *runs)
    default:
        fail("-mode must be sequential or parallel")
    }
}

func runSequential(a, b Matrix, runs int) {
    var rowsTotal time.Duration
    var columnsTotal time.Duration
    correct := true

    for run := 0; run < runs; run++ {
        rowsStarted := time.Now()
        byRows := multiplyByRows(a, b)
        rowsTotal += time.Since(rowsStarted)

        columnsStarted := time.Now()
        byColumns := multiplyByColumns(a, b)
        columnsTotal += time.Since(columnsStarted)

        correct = correct && matricesEqual(byRows, byColumns)
    }

    fmt.Printf("Matrix size: %d x %d\n", a.N, a.N)
    fmt.Printf("Runs: %d\n", runs)
    fmt.Printf("Average by rows:    %.6f s\n", average(rowsTotal, runs))
    fmt.Printf("Average by columns: %.6f s\n", average(columnsTotal, runs))
    fmt.Printf("Results are equal: %t\n", correct)
}

type workerStats struct {
    count   int
    total   time.Duration
    correct bool
}

func runParallel(a, b Matrix, workers []int, runs int) {
    stats := make([]workerStats, len(workers))
    for index, count := range workers {
        stats[index] = workerStats{count: count, correct: true}
    }

    var sequentialTotal time.Duration
    for run := 0; run < runs; run++ {
        referenceStarted := time.Now()
        reference := multiplyByRows(a, b)
        sequentialTotal += time.Since(referenceStarted)

        for index := range stats {
            started := time.Now()
            result, err := multiplyParallel(a, b, stats[index].count)
            stats[index].total += time.Since(started)
            if err != nil {
                fail(err.Error())
            }
            stats[index].correct = stats[index].correct && matricesEqual(reference, result)
        }
    }

    sequentialAverage := average(sequentialTotal, runs)
    fmt.Printf("Matrix size: %d x %d\n", a.N, a.N)
    fmt.Printf("Runs: %d\n", runs)
    fmt.Printf("Logical CPUs: %d\n", runtime.GOMAXPROCS(0))
    fmt.Printf("Average sequential: %.6f s\n\n", sequentialAverage)

    for _, stat := range stats {
        workerAverage := average(stat.total, runs)
        fmt.Printf(
            "Workers: %2d | Average: %.6f s | Speedup: %.2fx | Correct: %t\n",
            stat.count,
            workerAverage,
            sequentialAverage/workerAverage,
            stat.correct,
        )
    }
}

func average(total time.Duration, runs int) float64 {
    return total.Seconds() / float64(runs)
}

func parseWorkers(value string) ([]int, error) {
    parts := strings.Split(value, ",")
    workers := make([]int, 0, len(parts))
    for _, part := range parts {
        count, err := strconv.Atoi(strings.TrimSpace(part))
        if err != nil || count <= 0 {
            return nil, fmt.Errorf("invalid worker count %q", part)
        }
        workers = append(workers, count)
    }
    return workers, nil
}

func fail(message string) {
    fmt.Println("Error:", message)
    flag.Usage()
    os.Exit(2)
}
\end{lstlisting}

\begin{lstlisting}[caption={Файл 	exttt{matrix.go}}]
package main

import "math/rand"

type Matrix struct {
    N    int
    Data []int64
}

func newMatrix(n int) Matrix {
    return Matrix{N: n, Data: make([]int64, n*n)}
}

func (m Matrix) at(row, column int) int64 {
    return m.Data[row*m.N+column]
}

func (m Matrix) set(row, column int, value int64) {
    m.Data[row*m.N+column] = value
}

func generateMatrix(n, low, high int, seed int64) Matrix {
    if n <= 0 {
        panic("matrix size must be positive")
    }
    if low > high {
        panic("low bound must not exceed high bound")
    }

    rng := rand.New(rand.NewSource(seed))
    matrix := newMatrix(n)
    for index := range matrix.Data {
        matrix.Data[index] = int64(rng.Intn(high-low+1) + low)
    }
    return matrix
}

func matricesEqual(first, second Matrix) bool {
    if first.N != second.N || len(first.Data) != len(second.Data) {
        return false
    }
    for index, value := range first.Data {
        if value != second.Data[index] {
            return false
        }
    }
    return true
}
\end{lstlisting}

\begin{lstlisting}[caption={Файл 	exttt{multiply.go}}]
package main

import (
    "fmt"
    "math"
    "sync"
)

func multiplyByRows(a, b Matrix) Matrix {
    ensureCompatible(a, b)
    c := newMatrix(a.N)

    for row := 0; row < a.N; row++ {
        multiplyBlock(a, b, c, row, row+1, 0, a.N)
    }
    return c
}

func multiplyByColumns(a, b Matrix) Matrix {
    ensureCompatible(a, b)
    c := newMatrix(a.N)

    for column := 0; column < a.N; column++ {
        multiplyBlock(a, b, c, 0, a.N, column, column+1)
    }
    return c
}

func multiplyBlock(a, b, c Matrix, rowStart, rowEnd, columnStart, columnEnd int) {
    for row := rowStart; row < rowEnd; row++ {
        rowOffset := row * a.N
        for column := columnStart; column < columnEnd; column++ {
            var value int64
            for k := 0; k < a.N; k++ {
                value += a.Data[rowOffset+k] * b.Data[k*a.N+column]
            }
            c.Data[rowOffset+column] = value
        }
    }
}

type interval struct {
    start int
    end   int
}

func splitRange(size, parts int) []interval {
    ranges := make([]interval, parts)
    for part := range parts {
        ranges[part] = interval{
            start: part * size / parts,
            end:   (part + 1) * size / parts,
        }
    }
    return ranges
}

func chooseGrid(workers, n int) (rowBlocks, columnBlocks int, err error) {
    if workers <= 0 {
        return 0, 0, fmt.Errorf("worker count must be positive")
    }
    if workers > n*n {
        return 0, 0, fmt.Errorf("too many workers for a %dx%d matrix", n, n)
    }

    upper := int(math.Sqrt(float64(workers)))
    for rowBlocks = upper; rowBlocks >= 1; rowBlocks-- {
        if workers%rowBlocks != 0 {
            continue
        }
        columnBlocks = workers / rowBlocks
        if rowBlocks <= n && columnBlocks <= n {
            return rowBlocks, columnBlocks, nil
        }
    }
    return 0, 0, fmt.Errorf("cannot split %d workers into non-empty blocks for size %d", workers, n)
}

func multiplyParallel(a, b Matrix, workers int) (Matrix, error) {
    ensureCompatible(a, b)
    rowBlocks, columnBlocks, err := chooseGrid(workers, a.N)
    if err != nil {
        return Matrix{}, err
    }

    c := newMatrix(a.N)
    rowRanges := splitRange(a.N, rowBlocks)
    columnRanges := splitRange(a.N, columnBlocks)

    var group sync.WaitGroup
    group.Add(workers)
    for _, rows := range rowRanges {
        for _, columns := range columnRanges {
            go func(rows, columns interval) {
                defer group.Done()
                multiplyBlock(a, b, c, rows.start, rows.end, columns.start, columns.end)
            }(rows, columns)
        }
    }
    group.Wait()

    return c, nil
}

func ensureCompatible(a, b Matrix) {
    if a.N == 0 || a.N != b.N || len(a.Data) != a.N*a.N || len(b.Data) != b.N*b.N {
        panic("matrices must be non-empty squares of the same size")
    }
}
\end{lstlisting}

\begin{lstlisting}[language={},caption={Файл 	exttt{go.mod}}]
module bmstu/parallel-lab1

go 1.26.3
\end{lstlisting}

\end{document}
